\documentclass[UTF8]{ctexart}
\usepackage[a4paper,margin=2.6cm]{geometry}
\usepackage{amsmath, amssymb, amsthm}
\usepackage{graphicx}
\usepackage[colorlinks=true,linkcolor=blue]{hyperref}

\title{复变函数学习笔记：第6.1节 Gamma 函数}
\author{}
\date{}

\newtheorem{theorem}{定理}[section]
\newtheorem{lemma}[theorem]{引理}
\newtheorem{definition}[theorem]{定义}
\newtheorem{corollary}[theorem]{推论}
\newtheorem{example}{例}[section]
\newtheorem{remark}{注记}[section]
\numberwithin{equation}{section}

\newcommand{\RR}{\mathbb{R}}
\newcommand{\CC}{\mathbb{C}}

\begin{document}

\maketitle

本节系统研究 Gamma 函数 \(\Gamma(s)\) 的定义、解析延拓、互补公式、增长阶以及无穷乘积表示。Gamma 函数是最重要的非初等函数之一，它自然地出现在许多分析计算中，其倒数 \(1/\Gamma(s)\) 是以非正整数为零点的“最简”整函数。

% ========== 1.1 ==========
\section{定义与半平面解析性}

对于 \(\operatorname{Re}(s)>0\)，Gamma 函数由反常积分定义：
\begin{equation}\label{eq:GammaDef}
\Gamma(s) = \int_{0}^{\infty} e^{-t} t^{s-1}\, dt,
\end{equation}
其中 \(t^{s-1}=e^{(s-1)\log t}\) 取实数对数的主支。

\begin{theorem}[半平面解析性]\label{thm:GammaHalfPlane}
\(\Gamma(s)\) 在右半平面 \(\operatorname{Re}(s)>0\) 上全纯，且仍由 \eqref{eq:GammaDef} 给出。
\end{theorem}

\begin{proof}
分三步完成。

\textbf{1. 收敛性。} 设 \(s=\sigma+it\)，\(\sigma>0\)。注意到 \(|e^{-t}t^{s-1}|=e^{-t}t^{\sigma-1}\)。在 \(t\to0^{+}\) 处，\(t^{\sigma-1}\) 可积（因 \(\sigma>0\)）；在 \(t\to\infty\) 处，指数衰减 \(e^{-t}\) 压制任意幂次增长。故反常积分绝对收敛，\(\Gamma(s)\) 有定义。

\textbf{2. 构造全纯的截断函数列。} 对每个 \(\varepsilon>0\)，令
\[
F_{\varepsilon}(s)=\int_{\varepsilon}^{1/\varepsilon} e^{-t} t^{s-1}\, dt,\qquad s\in\CC.
\]
被积函数关于 \(s\) 全纯，且关于 \((s,t)\) 联合连续，故每个 \(F_{\varepsilon}\) 是整函数。

\textbf{3. 在竖带上一致收敛。} 任取正数 \(0<\delta<M<\infty\)，考虑竖带
\[
S_{\delta,M}=\{s\in\CC:\delta<\operatorname{Re}(s)<M\}.
\]
记 \(\sigma=\operatorname{Re}(s)\)，对 \(s\in S_{\delta,M}\) 有：
\[
\begin{aligned}
|\Gamma(s)-F_{\varepsilon}(s)|
&\le \int_{0}^{\varepsilon} e^{-t}t^{\sigma-1}\, dt + \int_{1/\varepsilon}^{\infty} e^{-t}t^{\sigma-1}\, dt\\
&\le \int_{0}^{\varepsilon} t^{\delta-1}\, dt + C\int_{1/\varepsilon}^{\infty} e^{-t/2}\, dt\\
&= \frac{\varepsilon^{\delta}}{\delta} + 2C e^{-1/(2\varepsilon)}.
\end{aligned}
\]
右端与 \(s\) 无关且当 \(\varepsilon\to0\) 时趋于 0，故 \(F_{\varepsilon}\to\Gamma\) 在 \(S_{\delta,M}\) 上一致。一致收敛的全纯函数列极限仍全纯，因此 \(\Gamma(s)\) 在 \(S_{\delta,M}\) 上全纯。由 \(\delta,M\) 的任意性，\(\Gamma(s)\) 在整个右半平面全纯。
\end{proof}

% ========== 1.2 ==========
\section{解析延拓到整个复平面}

利用分部积分可得 Gamma 函数的基本函数方程。

\begin{lemma}[函数方程]\label{lem:FuncEq}
若 \(\operatorname{Re}(s)>0\)，则
\[
\Gamma(s+1)=s\Gamma(s).
\]
特别地，\(\Gamma(n+1)=n!\) 对 \(n=0,1,2,\dots\) 成立。
\end{lemma}

\begin{proof}
对截断积分用微积分基本定理：
\[
\int_{\varepsilon}^{1/\varepsilon} \frac{d}{dt}\bigl(e^{-t}t^{s}\bigr)\, dt
= -\int_{\varepsilon}^{1/\varepsilon} e^{-t}t^{s}\, dt + s\int_{\varepsilon}^{1/\varepsilon} e^{-t}t^{s-1}\, dt.
\]
左边 \(=e^{-1/\varepsilon}(1/\varepsilon)^{s}-e^{-\varepsilon}\varepsilon^{s}\to0\) 当 \(\varepsilon\to0\)。取极限即得 \(\Gamma(s+1)=s\Gamma(s)\)。由 \(\Gamma(1)=\int_{0}^{\infty}e^{-t}dt=1\) 递推得 \(\Gamma(n+1)=n!\)。
\end{proof}

借助函数方程，可将 \(\Gamma(s)\) 逐步向左半平面延拓。

\begin{theorem}[解析延拓]\label{thm:GammaCont}
\(\Gamma(s)\) 可解析延拓为整个 \(\CC\) 上的亚纯函数，唯一奇点是单极点 \(s=0,-1,-2,\dots\)。在 \(s=-n\) 处的留数为 \((-1)^{n}/n!\)。
\end{theorem}

\begin{proof}
对任意正整数 \(m\)，在 \(\operatorname{Re}(s)>-m\) 上定义
\[
F_{m}(s)=\frac{\Gamma(s+m)}{(s+m-1)(s+m-2)\cdots s}.
\]
因 \(\Gamma(s+m)\) 在 \(\operatorname{Re}(s)>-m\) 上全纯，故 \(F_{m}\) 在该半平面亚纯，单极点为 \(s=0,-1,\dots,-m+1\)。由引理 \ref{lem:FuncEq} 知在 \(\operatorname{Re}(s)>0\) 上 \(F_{m}(s)=\Gamma(s)\)。由解析延拓的唯一性，\(\Gamma(s)\) 即为所有 \(F_{m}\) 在全平面上的共同延拓。
\end{proof}

% ========== 1.3 ==========
\section{余元公式（互补公式）}

Gamma 函数关于直线 \(\operatorname{Re}(s)=1/2\) 具有优美的对称性。

\begin{theorem}[余元公式]\label{thm:Complement}
对所有 \(s\in\CC\)，
\begin{equation}
\Gamma(s)\Gamma(1-s)=\frac{\pi}{\sin\pi s}.
\end{equation}
特别地，\(\Gamma(1/2)=\sqrt{\pi}\)。
\end{theorem}

\begin{proof}
只需对 \(0<s<1\) 证明等式，然后由解析延拓的唯一性（第二章定理 4.8）推广到全平面。因为等式两边都是 \(\CC\) 上的亚纯函数，且在区间 \((0,1)\) 上有极限点。

当 \(0<s<1\) 时，利用定义：
\[
\Gamma(s)\Gamma(1-s)=\int_{0}^{\infty}\int_{0}^{\infty} e^{-t-u} t^{s-1}u^{-s}\, dt\, du.
\]
作代换 \(u=vt\)（\(t>0\) 固定），得
\[
\Gamma(s)\Gamma(1-s)=\int_{0}^{\infty}\int_{0}^{\infty} e^{-t(1+v)} v^{-s}\, dv\, dt.
\]
交换积分次序，先积 \(t\)：
\[
\int_{0}^{\infty} e^{-t(1+v)}\, dt = \frac{1}{1+v},
\]
于是
\[
\Gamma(s)\Gamma(1-s)=\int_{0}^{\infty} \frac{v^{-s}}{1+v}\, dv.
\]
再令 \(v=e^{x}\)，化为
\[
\int_{-\infty}^{\infty} \frac{e^{(1-s)x}}{1+e^{x}}\, dx = \frac{\pi}{\sin\pi s},
\]
其中最后一步用到第三章例 2 的围道积分结果（矩形围道）。
\end{proof}

% ========== 1.4 ==========
\section{倒数 \(1/\Gamma(s)\) 及其增长阶}

从余元公式可直接读出倒数的结构。

\begin{theorem}[倒数 \(1/\Gamma(s)\)]\label{thm:RecipGamma}
\begin{enumerate}
\item \(1/\Gamma(s)\) 是整函数，只在 \(s=0,-1,-2,\dots\) 有单零点，其余点非零。
\item \(1/\Gamma(s)\) 的增长阶为 1，即对任意 \(\varepsilon>0\)，
\[
\left|\frac{1}{\Gamma(s)}\right|\le C(\varepsilon) e^{c|s|^{1+\varepsilon}}.
\]
\end{enumerate}
\end{theorem}

\begin{proof}
\textbf{(i)} 由余元公式 \(\displaystyle\frac{1}{\Gamma(s)}=\Gamma(1-s)\frac{\sin\pi s}{\pi}\)。\(\Gamma(1-s)\) 在 \(s=1,2,3,\dots\) 有单极点，而 \(\sin\pi s\) 恰在这些点有单零点，相互抵消；在 \(s=0,-1,-2,\dots\) 处 \(\Gamma(1-s)\) 全纯非零，\(\sin\pi s\) 有单零点。故 \(1/\Gamma(s)\) 为整函数且仅在非正整数点处有单零点。

\textbf{(ii)} 利用积分与级数的混合表示：
\[
\frac{1}{\Gamma(s)}= \underbrace{\left(\sum_{n=0}^{\infty}\frac{(-1)^{n}}{n!(n+1-s)}\right)\frac{\sin\pi s}{\pi}}_{I(s)}
+\underbrace{\left(\int_{1}^{\infty} e^{-t}t^{-s}\, dt\right)\frac{\sin\pi s}{\pi}}_{J(s)}.
\]
对 \(J(s)\)：由 \(|\int_{1}^{\infty} e^{-t}t^{-s}dt|\le e^{(|s|+1)\log(|s|+1)}\) 及 \(|\sin\pi s|\le e^{\pi|s|}\)，得 \(|J(s)|\le c_{1} e^{c_{2}|s|\log|s|}\)。

对 \(I(s)\)：当 \(|\operatorname{Im}(s)|\ge1\) 时，分母模有正下界，级数被 \(\sum 1/n!\) 控制；当 \(|\operatorname{Im}(s)|\le1\) 时，分离出最接近 \(s\) 的一项，利用 \(\sin\pi s\) 抵消极点，其余项一致有界。综上，两部分均满足所需上界。因此增长阶 \(\le1\)。

\noindent\textbf{增长阶恰为 1 的补充证明：}
取 \(s=-k-1/2\)（\(k\) 为正整数）。由余元公式，
\[
\left|\frac{1}{\Gamma(-k-1/2)}\right|
= \Gamma(k+3/2)\cdot\frac{|\sin\pi(-k-1/2)|}{\pi}
= \frac{1}{\pi}\Gamma(k+3/2) > \frac{1}{\pi} k!.
\]
由斯特林公式，\(k! \sim \sqrt{2\pi k}(k/e)^{k}\)，其主导项为 \(e^{k\log k}\)。若增长阶 \(\rho<1\)，则存在 \(\varepsilon>0\) 使 \(|1/\Gamma(s)|\le A e^{B|s|^{1-\varepsilon}}\)；但当 \(k\) 很大时，下界 \(e^{k\log k}\) 远快于上界 \(e^{B k^{1-\varepsilon}}\)，矛盾。因此增长阶不能小于 1。结合上界估计，增长阶恰为 1。
\end{proof}

% ========== 1.5 ==========
\section{Weierstrass 乘积公式}

利用 Hadamard 分解定理可得 \(1/\Gamma(s)\) 的无穷乘积表示，其中自然地引出 Euler 常数 \(\gamma\)。

\begin{definition}[Euler 常数]
Euler 常数定义为
\[
\gamma = \lim_{N\to\infty}\left(\sum_{n=1}^{N}\frac{1}{n} - \log N\right).
\]
该极限的存在性可以通过均值定理（将序列写成收敛级数）或单调有界原理证明。
\end{definition}

\begin{theorem}[Weierstrass 乘积公式]\label{thm:ProdFormula}
对所有 \(s\in\CC\)，
\[
\frac{1}{\Gamma(s)} = e^{\gamma s} s \prod_{n=1}^{\infty} \left(1+\frac{s}{n}\right) e^{-s/n}.
\]
\end{theorem}

\begin{proof}
由定理 \ref{thm:RecipGamma}，\(1/\Gamma(s)\) 是增长阶为 1 的整函数，零点为单零点 \(s=0,-1,-2,\dots\)。由 Hadamard 分解定理，存在一次多项式 \(P(s)=As+B\) 使得
\[
\frac{1}{\Gamma(s)} = e^{As+B} s \prod_{n=1}^{\infty} \left(1+\frac{s}{n}\right) e^{-s/n}.
\]
两边除以 \(s\) 并令 \(s\to0\)：左边 \(\to1\)（因 \(s\Gamma(s)\to1\)），右边 \(\to e^{B}\)，故 \(B=0\)。

令 \(s=1\)，利用 \(\Gamma(1)=1\) 得
\[
e^{-A} = \prod_{n=1}^{\infty} \left(1+\frac{1}{n}\right) e^{-1/n}.
\]
取部分积：
\[
\prod_{n=1}^{N} \left(1+\frac{1}{n}\right) e^{-1/n}
= (N+1) \exp\!\left(-\sum_{n=1}^{N}\frac{1}{n}\right)
= \frac{N+1}{N} e^{-(\sum_{n=1}^{N}1/n - \log N)} \to e^{-\gamma}.
\]
故 \(e^{-A}=e^{-\gamma}\)，即 \(A=\gamma\)。代入即得乘积公式。
\end{proof}

\end{document}